\begin{textsample}{POS Dim 3 – al – Score 15.00 – al/al\_inf\_18.txt}  \label{ex:f3_pos_219}
3.
DIREITOS DE APRENDIZAGENS E CAMPOS DE EXPERIÊNCIAS Conforme apresentados, a interação e a \textbf{brincadeira} são eixos estruturantes para a prática pedagógica na Educação \textbf{Infantil}, defendidos pela BNCC.
Nessa perspectiva, determina que seis direitos de aprendizagem e desenvolvimento sejam assegurados para as \textbf{crianças} nesta prática enquanto sujeitos sócio históricos e culturais.
Sendo assim, tais eixos apontam caminhos que assegurem a compreensão de que as \textbf{crianças} aprendem em situações nas quais possam desempenhar um papel ativo em ambientes que as \textbf{convidem} a vivenciar desafios e a sentirem -se provocadas a resolvê -los, nas quais possam construir significados sobre si, os outros e o mundo social e natural.
Compreendendo garantia de direitos para cumprimento de deveres em busca de uma sociedade alagoana justa, democrática e inclusiva, a BNCC aponta os direitos de aprendizagem, que são conviver, \textbf{brincar}, participar, explorar, expressar e conhecer -se, conforme descritos a seguir. 3.1 Dimensão conceitual dos Direitos de Aprendizagem : Conviver com outras \textbf{crianças} e \textbf{adultos}, em \textbf{pequenos} e grandes grupos, utilizando diferentes linguagens, ampliando o conhecimento de si e do outro, o respeito em relação à cultura e às diferenças entre as pessoas.
Nesse sentido, o direito de conviver é construído no estreitamento de vínculos afetivos, no compartilhamento de atitudes de solidariedade e no respeito aos tempos e espaços de cada um.
Conviver é no seu sentido mais amplo aprender a lidar emocionalmente com as diferenças, é aprender a respeitar a opinião do outro e abrir espaços para diálogos em situações de conflitos colaborando para um espaço de interação e convivência saudável. \textbf{Brincar} cotidianamente de diversas formas, em diferentes espaços e tempos, com diferentes \textbf{parceiros} ( \textbf{crianças} e \textbf{adultos} ), ampliando e diversificando seu acesso a produções culturais, seus conhecimentos, sua imaginação, sua criatividade, suas experiências emocionais, corporais, \textbf{sensoriais}, expressivas, cognitivas, sociais e \textbf{relacionais}.
No universo \textbf{infantil} tudo é possível, porque \textbf{brincar} é aguçar a imaginação é atribuir significado ao mundo.
O \textbf{brincar} vai muito além de um passatempo, de certo, é parte essencial para construção dos processos mentais do desenvolvimento \textbf{infantil}, pois garante a interação e potencializa o processo de aprendizagem.
Participar ativamente, com \textbf{adultos} e outras \textbf{crianças}, tanto do planejamento da gestão da escola e das atividades propostas pelo educador quanto da realização das atividades da vida cotidiana, tais como a escolha das \textbf{brincadeiras}, dos materiais e dos ambientes, desenvolvendo diferentes linguagens e elaborando conhecimentos, decidindo e se posicionando.
Explorar movimentos, \textbf{gestos}, \textbf{sons}, formas, \textbf{texturas}, cores, palavras, emoções, transformações, relacionamentos, histórias, objetos, elementos da natureza, na escola e fora dela, ampliando seus saberes sobre a cultura, em suas diversas modalidades : as artes, a escrita, a ciência e a tecnologia.
Expressar, como sujeito dialógico, criativo e sensível, suas necessidades, emoções, sentimentos, dúvidas, hipóteses, descobertas, opiniões, questionamentos, por meio de diferentes linguagens.
Conhecer -se e construir sua identidade pessoal, social e cultural, constituindo uma imagem positiva de si e de seus grupos de pertencimento, nas diversas experiências de \textbf{cuidados}, interações, \textbf{brincadeiras} e linguagens vivenciadas na \textbf{instituição} escolar e em seu contexto familiar e comunitário.

% matched lemmas: adulto, brincadeira, brincar, convidar, criança, cuidado, gesto, infantil, instituição, parceiro, pequeno, relacional, sensorial, som, textura
\end{textsample}
